\documentclass[10pt,a4paper]{article}

\usepackage[T1]{fontenc}
\usepackage[utf8]{inputenc}
\usepackage{lmodern}
\usepackage{amsmath,amssymb,amsthm}
\usepackage{booktabs,array}
\usepackage{geometry}
\usepackage{xcolor}
\usepackage{enumitem}
\usepackage{microtype}
\usepackage{hyperref}
\usepackage{fancyhdr}

\geometry{margin=1in}
\hypersetup{
  colorlinks=true,
  linkcolor=blue!55!black,
  citecolor=blue!55!black,
  urlcolor=blue!55!black,
  pdftitle={HPPK with Nothing: A Universal All-Zero Forgery against HPPK DS},
  pdfauthor={Martin Feussner},
  pdfsubject={Cryptanalytic audit report for IACR ePrint 2023/1768},
  pdfkeywords={digital signatures, HPPK, all-zero signature, universal forgery, cryptanalysis}
}
\setlist{nosep,leftmargin=*}
\setlength{\parindent}{1em}
\setlength{\parskip}{0.25em}
\pagestyle{fancy}
\fancyhf{}
\fancyhead[L]{HPPK with Nothing}
\fancyhead[R]{14 October 2026}
\fancyfoot[C]{\thepage}
\renewcommand{\headrulewidth}{0.3pt}

\newtheorem{theorem}{Theorem}
\newtheorem{lemma}{Lemma}
\newcommand{\Fp}{\mathbb{F}_p}

\title{\bfseries HPPK with Nothing:\\
A Universal All-Zero Forgery against HPPK DS}
\author{Martin Feussner\\
\small Selmer Center, University of Bergen\\
\small \texttt{martin.feussner@uib.no}}
\date{14 October 2026}

\begin{document}
\maketitle

\noindent\fcolorbox{black!50}{black!4}{\parbox{0.94\textwidth}{
This audit used the OpenAI model Daybreak Blue at the ultra reasoning level.
OpenAI Codex generated and selected this paper's title and intentionally chose
``HPPK with Nothing'' as wordplay: HPPK visually echoes ``happy,'' while the
forgery consists entirely of zeros. OpenAI Codex independently reconstructed,
derived, and verified the attack and performed AI-assisted attack-history
research, algebraic cryptanalysis, implementation, experimentation, and
technical writing. Separately tasked AI agents performed algebra/domain
review, hostile adversarial review, and fresh-context independent
reproduction, and a final exact-target history recheck. Primary reproduction,
adversarial review, independent reproduction, final history review, and
independent AI Publication Review passed. The result is classified as a
verified public design attack. No independent human verification was
performed; it is optional and is not a publication requirement.}}

\begin{abstract}
We describe a public-key-only universal forgery against the Homomorphic
Polynomial Public Key Cryptography for Quantum-secure Digital Signature
(HPPK DS) construction in IACR ePrint 2023/1768. For every public coefficient,
the verifier computes two Barrett-style expressions: each contains an outer
product with one signature component and a floor whose numerator contains that
same component. Substituting $F=H=0$ annihilates both outer products and sets
both floor arguments to zero, so the final verification equation is $0=0$ for
every message, structurally valid public key, and verifier-noise choice. The
advertised SHA-256, SHA-384, and
SHA-512 profiles split a complete hash into four field-sized representatives;
the full mathematical forgery is therefore four copies of $(0,0)$. It makes
zero signing queries and uses no secret information.

The proof covers all three rows of the current target's Table~1 and all fifteen
rows of the benchmark update's Table~2. A standard-library Python reproducer
checks the complete four-pair form and live controls for all eighteen rows. A
separate experiment against public \texttt{xtaci/hppk} commit
\texttt{a66c951eef2f231a228e8121a7a6391b991146e1} accepted 45 of 45
attacker-created, JSON-round-tripped zero objects across 15 fresh keys through
the ordinary exported verifier, with zero attack signing queries. That source
uses a different digest boundary, signature format, and parameter generation;
it is corroboration, not a row-exact implementation of the paper tables.
A fresh-context agent independently recreated the paper and source tests. Its
paper run passed 72/72 complete four-pair forgeries across all 18 rows, and its
pinned-source run accepted 192/192 attacker checks across 12 fresh keys and 48
fresh messages while 96/96 honest checks accepted and 288/288 negative
controls rejected. The attack used no signing queries or secret reads. Final
exact-target history review and independent Publication Review passed.
\end{abstract}

\section{Introduction}

Kuang, Perepechaenko, Sayed, and Lou propose HPPK DS by extending polynomial
public-key constructions over two hidden rings and transforming modular
multiplication into a public verification relation using Barrett-style
quotients~\cite{target}. The paper argues that the resulting floor terms make
forgery nonlinear and assigns exponential work to generic searches over the
signature values.

The boundary value zero bypasses that intended obstacle. Each transformed
coefficient contains an outer product with $F$ or $H$ and a Barrett correction
whose floor argument contains the same component. Both components may be zero
under the stated mathematical domain, and the verifier has no nonzero check.
The attack therefore requires no search: for every hash segment, output the
pair $(0,0)$.

This observation is elementary. We make no claim that zero substitution is a
new cryptanalytic technique. The contribution is the exact application to the
ordinary HPPK DS verifier, its complete four-segment signature, every
advertised row, and a separately scoped public-source implementation check.

\paragraph{Claim boundary.}
The paper-level claim is a universal forgery against the equations in the
current ePrint target and the benchmark's clarified row set. The paper defines
no canonical byte encoding. No public or reviewed row-exact implementation or
archived vector was located for a table row. The public Go source is tested
under its own source-specific format and parameters and is never presented as
execution of a paper row.

\section{Target, versions, and parameters}

\subsection{Pinned target}

The target is \emph{Homomorphic Polynomial Public Key Cryptography for
Quantum-secure Digital Signature}, IACR ePrint 2023/1768~\cite{target}. The
pinned current PDF has SHA-256 digest
\begin{center}
\small\ttfamily
faa261feaf9e1f030273d311bfe642504a82d9f0759524446f9c86afbbc269c1.
\end{center}
The official record reports receipt on 15 November 2023 and revision on
17 November 2023; the body matches arXiv 2311.08967v2. Later arXiv v3 is useful
lineage but is not silently substituted for the current ePrint target.

The performance and parameter update is \emph{Benchmark Performance of
Homomorphic Polynomial Public Key Cryptography for Key Encapsulation and
Digital Signature Schemes}, IACR ePrint 2024/019~\cite{benchmark}. Its pinned
PDF has SHA-256 digest
\begin{center}
\small\ttfamily
1ee259abede8b393f26e021fccbd41774b212146717a63f646aee282273198a7.
\end{center}

\subsection{Advertised rows}

The target Table~1 gives three security levels but omits a concrete prime and
an explicit $(n,\lambda,m)$ tuple. Its surrounding construction supports the
contextual values $n=\lambda=1,m=2$; we do not turn that inference into an
exact table cell. Table~\ref{tab:targetrows} records the values that are
actually printed.

\begin{table}[ht]
\centering
\small
\begin{tabular}{@{}lrrrrl@{}}
\toprule
Level & field bits & hidden bits & Barrett bits & signature bytes & hash\\
\midrule
I   & 64  & 128 & 160 & 128 & SHA-256\\
III & 96  & 192 & 224 & 192 & SHA-384\\
V   & 128 & 256 & 288 & 256 & SHA-512\\
\bottomrule
\end{tabular}
\caption{Current ePrint 2023/1768 Table~1 rows. A complete hash contains four
field-sized segments in every row.}
\label{tab:targetrows}
\end{table}

Benchmark Table~2 makes fifteen tuples explicit. All use
$n=\lambda=1$; Table~\ref{tab:benchrows} groups the five rows at each level.
The package's \texttt{affected-rows.json} retains every row individually.

\begin{table}[ht]
\centering
\scriptsize
\begin{tabular}{@{}lllll@{}}
\toprule
Level & $p$ & $(m,L,K)$ tuples & hash & signature bytes\\
\midrule
I & $2^{64}-59$ & $(2,144,288)$, $(2,144,208)$, $(2,144,176)$,
  $(1,144,208)$, $(1,144,176)$ & SHA-256 & 144\\
III & $2^{96}-17$ & $(2,208,416)$, $(2,208,272)$, $(2,208,240)$,
  $(1,208,272)$, $(1,208,240)$ & SHA-384 & 208\\
V & $2^{128}-159$ & $(2,272,544)$, $(2,272,336)$, $(2,272,304)$,
  $(1,272,336)$, $(1,272,304)$ & SHA-512 & 272\\
\bottomrule
\end{tabular}
\caption{All fifteen ePrint 2024/019 Table~2 rows.}
\label{tab:benchrows}
\end{table}

\subsection{Claimed security}

The original HPPK DS paper sizes its parameter sets from the claimed
$O(S_1S_2/p)=O(p^3)$ key-recovery cost and reports estimated entropies of 192,
288, and 384 bits in Table~1. Propositions~3--5 separately present generic
forgery searches with claimed cost $O(S_1S_2)$. The companion benchmark later
presents an optimized modulus-recovery search with claimed total cost
$O(2^L)$; its Table~2 sets $L=2\lvert p\rvert_2+16$ and labels Levels I, III,
and V with 144, 208, and 272 bits. The all-zero forgery bypasses these searches:
it constructs the fixed four-pair signature directly and does not depend on
either claimed entropy column.

\section{HPPK DS verification}

Let $p$ be the public field modulus. The scheme forms product polynomials
$P=fB$ and $Q=hB$, encrypts their coefficients over hidden rings with moduli
$S_1,S_2$, and publishes Barrett-transformed data. For coefficient indices
$i\in\{0,\ldots,n+\lambda\}$ and $j\in\{1,\ldots,m\}$, write
\[
 p'_{ij}=\beta P_{ij}\bmod p,\qquad
 q'_{ij}=\beta Q_{ij}\bmod p,
\]
\[
 s_1=\beta S_1\bmod p,\qquad s_2=\beta S_2\bmod p,
\]
and let $\mu_{ij},\nu_{ij}$ be the published integer quotient
approximations for a positive power of two $R$.

For one field-sized message representative $x$, signing samples
$\alpha\in\Fp$ and produces residues $F$ and $H$ modulo the two hidden rings.
The stated domain includes $\alpha=0$, and no signing or verification rule
excludes the resulting zero pair.

Target Equation~(28) embeds the signature in the two coefficient equations
\begin{align}
 U_{ij}(H) &= H p'_{ij}
  -s_1\left\lfloor\frac{H\mu_{ij}}{R}\right\rfloor \bmod p,
  \label{eq:u-coefficient}\\
 V_{ij}(F) &= F q'_{ij}
  -s_2\left\lfloor\frac{F\nu_{ij}}{R}\right\rfloor \bmod p.
  \label{eq:v-coefficient}
\end{align}
The target aggregates these coefficients with powers $x^i$ and verifier noise
values $u_j$, then accepts when the two aggregates are equal. The benchmark
checks the corresponding equality separately for each $j$. Neither form tests
that $F$ or $H$ is nonzero.

Hashes wider than the field are split into field-sized segments and each
segment is signed independently. SHA-256 over a 64-bit field, SHA-384 over a
96-bit field, and SHA-512 over a 128-bit field each produce exactly four
segments. Thus a full mathematical signature contains four pairs.

\section{Universal all-zero forgery}

\begin{lemma}[Zero coefficients]
For every structurally valid public coefficient, field representative, and
positive Barrett power $R$, substituting $F=H=0$ in
Equations~\eqref{eq:u-coefficient} and~\eqref{eq:v-coefficient} gives
$U_{ij}=V_{ij}=0$.
\end{lemma}

\begin{proof}
Both products outside the floor vanish. Inside the floor,
$0\mu_{ij}/R=0$ and $0\nu_{ij}/R=0$, so both floor values are zero. Reduction
modulo $p$ preserves zero. This is exact and does not depend on the quality of
the Barrett approximation.
\end{proof}

\begin{theorem}[Complete-hash universal forgery]
For each advertised HPPK DS row, the four-pair signature
\[
 \bigl((0,0),(0,0),(0,0),(0,0)\bigr)
\]
is accepted for every complete message hash, every structurally valid public
key, and every verifier-noise choice by the ordinary published verifier.
\end{theorem}

\begin{proof}
Apply the lemma to every coefficient for one segment. Every polynomial sum is
zero, so both the target's noise-weighted equation and every benchmark
per-$j$ equation compare zero with zero. Repeat independently for all four
hash segments. All four comparisons accept. The conclusion is independent of
the message representatives and public coefficients.
\end{proof}

The attacker chooses any fresh message, emits eight zero integers, and makes
no signing query. Work and output are linear only in the fixed four-segment
signature length. Under any natural fixed-width unsigned encoding, the object
would be an all-zero byte string of the printed signature length. Because the
papers do not define a canonical wire format, this observation is illustrative
rather than an interoperability claim.

\section{Paper-level executable reproduction}

The public reproducer uses only the Python standard library. It transcribes
all three target rows and all fifteen benchmark rows. For each row it hashes
explicit messages with the named hash function, splits the digest into four
field-sized representatives, constructs four zero pairs, evaluates every
coefficient with integer floor division, and checks both the per-$j$ and
noise-aggregated equalities.

Target Table~1 does not publish a prime. For those three records the executable
uses the corresponding later benchmark prime as a labeled arithmetic fixture.
The result does not depend on that choice; it is not reported as a target-row
parameter. The benchmark records use their exact Table~2 primes, dimensions,
$L$, and $K$. The coefficients are deterministic audit fixtures, not author
keys. Accordingly, the executable checks the exact universal identity while
making no row-exact implementation claim.

The quick reference checks one complete hash per row, for 18 messages and 72
zero pairs. The full reference checks four complete hashes per row, for 72
messages and 288 zero pairs. Both pass every row with zero signing queries.
Each row also requires:
\begin{itemize}
  \item rejection of three-pair and five-pair objects;
  \item rejection when only $F$ or only $H$ is changed to one under a fixture
        whose constant coefficient makes the comparison path live;
  \item rejection of the zero object under a test-only per-segment repair; and
  \item acceptance of a balanced nonzero equation fixture under that repair.
\end{itemize}
The last item is a path control, not a claim that the fixture is produced by
the unavailable row-exact signer.

\section{Separate public Go implementation evidence}

The public \texttt{xtaci/hppk} repository was pinned at commit~\cite{source}
\begin{center}
\small\ttfamily
a66c951eef2f231a228e8121a7a6391b991146e1.
\end{center}
The retained archive has SHA-256 digest
\begin{center}
\small\ttfamily
f013a87fab81955e814d23e0ad869454523898b560285dc9664a1b474a61f28c.
\end{center}
Its exact 19-file inventory validated before and after the experiment, and
\texttt{hppk.go} remained SHA-256
\begin{center}
\small\ttfamily
052291b33de3c5a6a64d5de1752da6c3af01b92cdbbdf3862545c60dee1dea91.
\end{center}

This source is materially different from the papers. It consumes one entire
digest as one integer, requires key order at least five, generates different
hidden moduli and Barrett values, and places
\texttt{Beta}, \texttt{S1Verify}, \texttt{S2Verify}, \texttt{U},
\texttt{V}, and \texttt{K} inside the signature object. No source
configuration is an exact advertised row.

The primary harness nevertheless tests the same zero boundary through the
ordinary exported \texttt{hppk.VerifySignature} function. For each trial it
uses a fresh source-generated key, builds a complete attacker object from
public array lengths and constants, serializes it to JSON, decodes it into a
fresh signature, and verifies the forged object before any signing call. Five
source-specific profiles were exercised:
\begin{itemize}
  \item custom primes $2^{64}-59$, $2^{96}-17$, and $2^{128}-159$, each at
        order five with SHA-256, SHA-384, and SHA-512 respectively; and
  \item the source's 2,049-bit \texttt{DefaultPrime} at orders five and ten.
\end{itemize}

Three fresh keys per profile and three distinct messages per key gave 15
unique public-key fingerprints and 45/45 accepted forgeries. The attack made
zero signing queries. An object with nonzero auxiliary values but $F=H=0$
also accepted for every key, isolating the two paper components as the cause.

After attack construction and verification, fifteen control-only calls to
\texttt{Sign} established honest acceptance. Honest changed-message,
$F+1$ mutation, targeted nonzero, malformed-length, nil-signature,
empty-digest, and wrong-key controls rejected as registered. These later
signing calls are never counted as attack queries.

All generated keys had the intended shape, and ordinary honest signatures
verified. The diagnostic \texttt{PublicKey.IsValid()} nevertheless returned
false. Key generation maps
coefficients into the hidden-ring range, whereas that diagnostic requires
them to be below the public prime. The ordinary verifier does not call
\texttt{IsValid}. This is a separate source consistency issue, not the attack
mechanism.

The host's Go 1.22.2 rejected the source's Go 1.22.5 module directive. A
prospectively registered fallback copied the verified source byte for byte to
an offline temporary GOPATH, disabled module and network resolution, and built
without source edits. Pre- and post-build recursive comparisons were empty.
The complete registered workflow used 24.59 seconds wall time, 19.20 CPU
seconds, 230,900 KiB peak resident memory, and one worker. The implementation
harness itself used 0.04 seconds wall time, 0.04 CPU seconds, and 6,144 KiB
peak resident memory.

\subsection{Fresh-context independent reproduction}

A separately tasked AI agent reconstructed the equation and source harnesses
from the public target materials and pinned source. At the paper level, all 3
target and 15 benchmark rows passed: 72/72 complete fresh-message signatures,
288/288 segments, and 1,440 coefficient pairs. Each complete signature used
four zero pairs. Signing queries and secret reads were zero.

The separate source run covered custom 64-, 96-, and 128-bit primes and the
source \texttt{DefaultPrime} at orders 5, 8, and 10. It generated 12 unique
fresh keys and 48 unique fresh messages. The ordinary exported verifier
accepted both all-zero and live-auxiliary attacker objects under correct and
wrong keys, for 192/192 attacker checks. The live objects retained $F=H=0$
while using nonzero in-range auxiliary values, legitimate lengths, and the
source-derived normal positive \texttt{K}. Honest direct and serialized checks
accepted 96/96; six negative-control classes rejected 288/288. Source
inventories matched exactly before and after. The source run used one worker,
0.34 seconds wall time, and 7,680 KiB peak resident memory.

Attempt 1 completed the paper checks, then stopped at compilation before a
source binary was created or run because of a pointer-slice bookkeeping type
error. A prospectively sealed one-line type correction enabled the sole source
execution. No source-result tuning or retry occurred.

\section{Specification inconsistencies and limitations}

The result survives several specification ambiguities, but the public report
must retain them:
\begin{enumerate}
  \item \textbf{No paper wire format.} The papers provide size formulas but no
        canonical component ranges, endianness, integer encoding, or parser.
        The complete paper claim is therefore four mathematical pairs.
  \item \textbf{No located row-exact executable.} No public or reviewed
        row-exact implementation or archived vector was located for any target
        or benchmark row. The proof is full-parameter because both coefficient
        expressions evaluate to zero independently of the parameter values,
        not because a table implementation was run.
  \item \textbf{Broken benchmark pseudocode.} Algorithm~5 duplicates $F$ and
        omits the evident $H$ line. Algorithm~6 prints an extra coefficient
        index and uses $s_1$ where the equations require $s_2$. Nearby prose
        gives conflicting signature sizes, and a caption gives Level-I
        $L=168$ against Table~2's $L=144$. At $F=H=0$, both printed floor
        arguments and both outer products are zero under either branch; the
        claim is anchored to target Equation~(28), its signing and verification
        text, benchmark Equation~(8), and Table~2.
  \item \textbf{Source divergence.} The one-object Go verifier result is
        source-specific corroboration and cannot be used as a four-segment or
        row-exact execution claim.
  \item \textbf{Attack type.} This is a universal fresh-message forgery, not
        secret-key or hidden-ring recovery.
  \item \textbf{Publication gate.} Independent Publication Review passed.
\end{enumerate}

\section{Attack history and attribution}

In public searches completed through 14 October 2026 at 14:47:51 UTC, we
located no previous report of the four-pair all-zero universal forgery, or of
another practical signing compromise, against the exact Equation~(28) HPPK DS
paper construction in ePrint 2023/1768. This is a bounded, dated conclusion
with moderate-high confidence. It cannot exclude nonpublic, deleted,
unpublished, or unindexed disclosures, and search and citation indexes can
lag. We claim neither absolute priority nor technique novelty.

The final search did locate a September 2026 public-key-only forgery against a
materially divergent HPPK-labelled implementation, reported in a draft pull
request and explicit public reproduction commit~\cite{hppk2pr}. That
implementation's attacker-supplied signature contains $(F,H,U,V)$. Its test
chooses $F=H=1$, computes $U$ and $V$ from the message and public key at
$p=65537$, and passes a verifier that separately checks
$U=\alpha H\bmod p$ and $V=\beta F\bmod p$. The paper target instead supplies
only $(F,H)$: its verifier derives every $U_{ij}(H)$ and $V_{ij}(F)$ internally
and compares the two polynomial sides under Equation~(28). The September
finding is a real prior implementation-specific forgery, but it neither
reports the four-pair all-zero attack nor executes an advertised paper row.
The repository owner is not a target-paper author, and no target association
was located. We classify it as related work against a divergent implementation,
not an attack on the exact paper construction.

The target paper's Propositions~3--5 present generic forgery searches with
claimed cost $O(S_1S_2)$ and do not give a practical advertised-parameter
compromise. Guo's
practical cryptanalysis of MPPK/DS attacks a distinct predecessor with a
different public structure and linear verifier~\cite{guo}. The September 2026
D-HPPK KEM preprint reports a Lattice Reconstitution Attack against the
original unmasked HPPK KEM~\cite{dhppk}; it does not analyze the ePrint
2023/1768 signature verifier, recover a DS signing witness, or demonstrate
fresh DS signing. These items are relevant history, but each falls outside the
exact Equation~(28) signature target.

\section{Narrow repair direction}

For every hash segment, verification should reject when both $F=0$ and $H=0$.
The signer must resample $\alpha$, or the resulting pair, whenever it would
emit the newly forbidden value. These checks block the exact paper-level
forgery in this report.

This is a narrow mitigation, not a proof of unforgeability. A complete revision
should define canonical component ranges and encoding and should repeat its
security and full-parameter analysis. The divergent Go format additionally
needs coherent range/domain validation and binding of its auxiliary values;
those questions are separate from the minimal paper repair.

\section{Reproducibility and review status}

The public artifact includes the full standard-library source, all 18 affected
rows, deterministic quick and full reference outputs, package checks, build
files, and this \LaTeX{} source. Ordinary runs use operating-system entropy in
their message labels; reference runs use a disclosed fixed fixture. Every run
is sequential with one worker and emits machine-readable JSON ending in
\texttt{ATTACK REPRODUCTION: PASS} on success.

Primary reproduction, hostile adversarial review, fresh-context independent
reproduction, final exact-target history review, and independent Publication
Review passed. The package is classified as a verified public design attack.
Its final file inventory is sealed by the accompanying \texttt{SHA256SUMS}
manifest.

\section*{Acknowledgements}
The computations were performed on the Norwegian Research and Education Cloud (NREC), using resources provided by the University of Bergen and the University of Oslo.

\begin{thebibliography}{9}

\bibitem{target}
Randy Kuang, Maria Perepechaenko, Mahmoud Sayed, and Dafu Lou.
\newblock Homomorphic Polynomial Public Key Cryptography for Quantum-secure
Digital Signature.
\newblock IACR Cryptology ePrint Archive, Report 2023/1768, 2023.
\newblock \url{https://eprint.iacr.org/2023/1768}.

\bibitem{benchmark}
Randy Kuang, Maria Perepechaenko, Dafu Lou, and Brinda Tank.
\newblock Benchmark Performance of Homomorphic Polynomial Public Key
Cryptography for Key Encapsulation and Digital Signature Schemes.
\newblock IACR Cryptology ePrint Archive, Report 2024/019, 2024.
\newblock \url{https://eprint.iacr.org/2024/019}.

\bibitem{guo}
Hao Guo.
\newblock Cryptanalysis of MPPK/DS.
\newblock IACR Cryptology ePrint Archive, Report 2023/453, 2023.
\newblock \url{https://eprint.iacr.org/2023/453}.

\bibitem{hppk2pr}
sergelen02.
\newblock Test: add HPPK DS security regression suite.
\newblock Draft GitHub pull request 1, created 9 September 2026.
\newblock \url{https://github.com/sergelen02/HPPK_2/pull/1}.
\newblock Public reproduction commit
\texttt{782f51815b011dbd16ac4dd5de48757e0fb69bf4}.

\bibitem{dhppk}
D-HPPK KEM preprint.
\newblock D-HPPK KEM: A Defactorized Homomorphic Polynomial Public Key
Encapsulation Mechanism for IND-CCA2 Security.
\newblock Preprints.org manuscript 202608.1587, version 3, September 2026.
\newblock \url{https://www.preprints.org/manuscript/202608.1587}.

\bibitem{source}
xtaci.
\newblock hppk public source repository.
\newblock Commit
\texttt{a66c951eef2f231a228e8121a7a6391b991146e1}.
\newblock \url{https://github.com/xtaci/hppk}.

\end{thebibliography}

\end{document}
